\documentclass[../../Bayesian_Inference.tex]{subfiles}

\begin{document}

\chapter{Fundamentals of Bayesian Inference} 

Consider an observed dataset $\mathcal{D}$ and a statistical model parameterized by an unknown parameter vector $\bm{\theta} \in \Theta$, where $\Theta$ denotes the parameter space. In the Bayesian framework, uncertainty about $\bm{\theta}$ before observing the data is represented by a prior distribution $p(\bm{\theta})$.

Given a particular value of $\bm{\theta}$, the observed data are characterized by the conditional distribution
\[
    p(\mathcal{D}\mid\bm{\theta}),
\]
which, when viewed as a function of $\bm{\theta}$ for fixed $\mathcal{D}$, defines the likelihood function.

The joint distribution of $\bm{\theta}$ and $\mathcal{D}$ can be factorized in two equivalent ways:
\[
    p(\bm{\theta},\mathcal{D})
    =
    p(\mathcal{D}\mid\bm{\theta})p(\bm{\theta}),
    \qquad 
    p(\bm{\theta},\mathcal{D})
    =
    p(\bm{\theta}\mid\mathcal{D})p(\mathcal{D}).
\]
Equating the two factorizations yields
\[
    p(\bm{\theta}\mid\mathcal{D})p(\mathcal{D})
    =
    p(\mathcal{D}\mid\bm{\theta})p(\bm{\theta}),
\]
and therefore Bayes' theorem gives
\begin{equation}
    \boxed{
    p(\bm{\theta}\mid\mathcal{D})
    =
    \frac{
        p(\mathcal{D}\mid\bm{\theta})p(\bm{\theta})
    }{
        p(\mathcal{D})
    }}
    .
\end{equation}

The denominator is obtained by marginalizing the joint distribution over the parameter space:
\begin{equation}
    \boxed{p(\mathcal{D})
    = 
    \int_{\Theta}
    p(\mathcal{D}\mid\bm{\theta})
    p(\bm{\theta})
    \mathrm{d}\bm{\theta}}
    .    
\end{equation} 

The quantity $p(\mathcal{D})$ is referred to as the marginal likelihood, or model evidence. It normalizes the posterior distribution so that
\[
    \int_{\Theta}
    p(\bm{\theta}\mid\mathcal{D})
    \mathrm{d}\bm{\theta}
    =
    1.
\]

For a fixed observed dataset $\mathcal{D}$, the marginal likelihood $p(\mathcal{D})$ does not depend on $\bm{\theta}$. Hence, the posterior distribution can be written up to a proportionality constant as
\begin{equation}
    \boxed{
    p(\bm{\theta}\mid\mathcal{D})
    \propto
    p(\mathcal{D}\mid\bm{\theta})
    p(\bm{\theta})}.
\end{equation}

Thus, the basic principle of Bayesian inference can be summarized as
\[
    \boxed{
    \text{Posterior}
    \propto
    \text{Likelihood}
    \times
    \text{Prior}}.
\]

Bayesian inference can also be used to quantify uncertainty about future observations. Let $\mathcal{D}^{*}$ denote a future observation or a set of future observations. By marginalizing over the uncertainty in $\bm{\theta}$, the posterior predictive distribution is
\[
    p(\mathcal{D}^{*}\mid\mathcal{D})
    =
    \int_{\Theta}
    p(\mathcal{D}^{*}\mid\bm{\theta},\mathcal{D})
    p(\bm{\theta}\mid\mathcal{D})
    \mathrm{d}\bm{\theta}.
\]

If $\mathcal{D}^{*}$ is conditionally independent of $\mathcal{D}$ given $\bm{\theta}$, i.e., $\mathcal{D}^{*} \perp \mathcal{D} \mid \bm{\theta}$, then
\[
    p(\mathcal{D}^{*}\mid\bm{\theta},\mathcal{D})
    =
    p(\mathcal{D}^{*}\mid\bm{\theta}),
\]
and the posterior predictive distribution reduces to
\begin{equation}
    \boxed{
    p(\mathcal{D}^{*}\mid\mathcal{D})
    =
    \int_{\Theta}
    p(\mathcal{D}^{*}\mid\bm{\theta})
    p(\bm{\theta}\mid\mathcal{D})
    \mathrm{d}\bm{\theta}}
    .
\end{equation} 
The posterior predictive distribution therefore propagates uncertainty in model the parameters into uncertainty about future observations by averaging the predictive distribution over the posterior distribution of $\bm{\theta}$. 


\end{document}